% !PRINTED 567
\weberpage{567}

%<<CONT
from which:
\weq{7}{a\partial - bc = \dfrac{x^2 - \Delta y^2}{4} = \mu\mu' = m,}
if
\weq{8}{\mu' = \dfrac{x - y\sqrt{\Delta}}{2}}
is the number of the field $\Omega$ conjugate to $\mu$.

If we take $y$ different from zero, then from (1), by multiplication by $y$, there follows:
\[ b\omega_2^2 + (a-\partial)\omega_2\omega_1 - c\omega_1^2 = 0, \]
and for this one can also write:
\[ (a\omega_1 + b\omega_2)\omega_2 = (c\omega_1 + \partial\omega_2)\omega_1 \]
or
\weq{9}{\omega = \dfrac{c + \partial\omega}{a + b\omega}.}
By (4) and (5) there follows from (6):
\weq{10}{\begin{aligned} a + b\omega &= \dfrac{x + (2A\omega+B)y}{2} = \mu, \\ c + \partial\omega &= \mu\omega, \end{aligned}}
or in homogeneous form:
\weq{11}{\begin{aligned} \mu\omega_1 &= a\omega_1 + b\omega_2, \\ \mu\omega_2 &= c\omega_1 + \partial\omega_2, \end{aligned}}
and by solving these linear equations:
\weq{12}{\begin{aligned} \mu'\omega_1 &= \phantom{-}\partial\omega_1 - b\omega_2, \\ \mu'\omega_2 &= -c\omega_1 + a\omega_2. \end{aligned}}

From this it follows that $\wp(\mu u)$ is a doubly-periodic function with the periods $\omega_1$, $\omega_2$, and since it is moreover an even function of $u$, it can be represented rationally by $\wp(u)$ (§ 21). If, then, we denote by $R$ and $P$ integral rational functions of $\wp(u)$ without common divisor, then
\weq{13}{\wp(\mu u) = \dfrac{R}{P}.}

Since the quotient $\wp(\mu u):\wp(u)$ remains finite for $u=0$, i.\,e. for $\wp(u)=\infty$, the degree of $R$ is higher by one unit than the degree of the denominator.
